En tres nodos en \emph{localhost}, con Python~3.13 y una conexión o un canal por cliente, sockets TCP con JSON tuvo menor latencia que gRPC con Protocol Buffers en las cuatro condiciones que medimos. Con 100~bytes y un cliente, la mediana fue de 0,045~ms frente a 0,291~ms, y sockets atendió unos 7700 mensajes por segundo frente a 2600 de gRPC. H1 se cumple: la diferencia está en cada repetición y los rangos de los dos mecanismos no se tocan.

El tamaño del mensaje redujo la ventaja relativa de sockets, como decía H2. Con 10~KB, gRPC tardó unas 4,2 veces más que sockets, frente a 6,5 veces con 100~bytes. Eso pasó porque sockets se hizo más lento con el mensaje grande y gRPC casi no cambió. La diferencia en milisegundos se quedó cerca de 0,23~ms. En estos tamaños, lo que más cuesta en gRPC es el trabajo fijo de cada llamada, no el contenido.

H3 no se confirmó. Con tres clientes a la vez, las medianas de los dos mecanismos subieron en proporción parecida, y no se puede decir que la ventaja de sockets haya crecido. gRPC aprovechó mejor la concurrencia en mensajes por segundo (los duplicó, mientras que sockets los subió un 59\,\%), aunque sockets siguió por encima.

gRPC fue más estable. Su envío más lento tardó 7,8~ms, mientras que sockets tuvo 13 envíos de más de 10~ms, cuya causa no identificamos. Para este caso, sockets es la opción con menos latencia; si se necesita que el tiempo de respuesta sea predecible, los datos favorecen a gRPC. La facilidad de desarrollo de cada mecanismo se trata en la sección de discusión y no cambia estos números.

Los relojes de Lamport funcionaron con los dos mecanismos: en los 20\,400 mensajes, el reloj de la recepción fue mayor que el del envío. La convergencia necesitó 2, 6 y 8 intercambios con 2, 3 y 4 nodos. El número crece con los nodos, pero no da los valores que pusimos en H4, porque la ronda final usa dos fichas y la hipótesis suponía una. En el tiempo, H4 no se confirmó porque algunas corridas incluyen medio segundo de arranque de los procesos.

Todo esto se midió en una sola máquina y sin red real. Con una red de verdad, el tiempo de transmisión se sumaría a los dos mecanismos, y esperamos que la diferencia relativa sea menor. Para repetir el experimento conviene que los nodos esperen a que todos los servidores estén escuchando antes del primer envío, y buscar la causa de los envíos lentos de sockets.

\subsection{Conclusiones individuales}

% Cada integrante escribe aquí su propia conclusión (mínimo 150 palabras):
% qué aprendió, qué le costó, qué haría distinto. Borrar la línea "Pendiente" al escribirla.

\subsubsection*{Castro Coello Mario Sebastián}
\textit{Pendiente.}

\subsubsection*{Loor Suárez Steffanny del Rocío}
Durante esta práctica pude comprender mejor cómo funciona la
comunicación entre procesos mediante sockets TCP y cómo
se relaciona con los sistemas distribuidos.
Mi participación en esta tarea se centró en desarrollar el servidor y
el cliente TCP, comprobar el intercambio de mensajes e
integrar esta comunicación con el reloj lógico de Lamport.
Con ello entendí que no basta con enviar información,
sino que también es necesario validar los datos y confirmar su
recepción y controlar las conexiones.

Una de las dificultades que encontré fue conseguir que los
procesos se comunicaran correctamente y manejar los errores
que podían aparecer cuando un nodo todavía no estaba disponible.
También tuve que revisar el funcionamiento de los mensajes de
confirmación y comprobar que la integración con el reloj de Lamport
mantuviera el orden lógico de los eventos.

Puedo dar como conclusión de esta práctica que según los resultados
TCP puede ofrecer una latencia menor que gRPC en las condiciones
evaluadas, aunque eso no significa que siempre sea la mejor alternativa.
Si volviera a realizar esta práctica, organizaría desde el inicio pruebas
más específicas para detectar los errores de conexión y verificaría
que todos los servidores estuvieran preparados antes de comenzar los intercambios.

\subsubsection*{Quintero Rico Luis David}
\textit{Pendiente.}

\subsubsection*{Rivas Piloso Dayver Yahir}
\textit{Pendiente.}
